% =====================================================================
%  Pacotes e configuracoes gerais do preambulo.
%  Adicione aqui os pacotes que o seu trabalho precisar.
% =====================================================================

% --- Codificacao e idioma --------------------------------------------
\usepackage[T1]{fontenc}
\usepackage[utf8]{inputenc}
\usepackage{lmodern}

% --- Elementos graficos ----------------------------------------------
\usepackage{graphicx}
\graphicspath{{figures/}}
\usepackage{float}
\usepackage{booktabs}
\usepackage{multirow}
\usepackage{longtable}

% --- Matematica ------------------------------------------------------
\usepackage{amsmath}
\usepackage{amssymb}

% --- Tipografia ------------------------------------------------------
\usepackage{microtype}
\usepackage{indentfirst}
\usepackage{color}
\usepackage{xcolor}

% --- Links -----------------------------------------------------------
% A abntex2 ja carrega o hyperref; aqui so ajustamos a aparencia.
% Para a versao de impressao, troque colorlinks por false.
\definecolor{lgLink}{RGB}{0,80,160}
\hypersetup{
  colorlinks=true,
  linkcolor=lgLink,
  citecolor=lgLink,
  urlcolor=lgLink,
  bookmarksdepth=4
}

% --- Espacamento -----------------------------------------------------
\setlength{\parindent}{1.3cm}
\setlength{\parskip}{0.2cm}
\OnehalfSpacing

% --- Comandos de conveniencia ----------------------------------------
% A abntex2 ja fornece \fonte{...} para a fonte de figuras e tabelas,
% exigida pela NBR 14724 logo abaixo do elemento. Use assim:
%   \fonte{Elaborado pelo autor (2026).}
% Adicione aqui os seus proprios comandos.
